%=============================================================================
% Section 12: Conclusion
%=============================================================================
\section{Conclusion}

The emergence of AI agents requires a dedicated architectural boundary for capability execution. This paper introduces the AI Execution Layer as a minimal runtime architecture for agent-driven systems.

The architecture separates cognition, orchestration, and execution through the logical model $S = (A, P, E)$. Within this boundary, the runtime is constrained to six execution responsibilities:
\begin{enumerate}
    \item Capability invocation
    \item Protocol translation
    \item Result normalization
    \item Connection handling
    \item Error propagation
    \item Session transport
\end{enumerate}

The theoretical foundation is defined by four principles:
\begin{enumerate}
    \item \textbf{Capability Abstraction:} Agents depend on capabilities, not providers
    \item \textbf{Cognition--Action Separation:} Intelligence and execution are separate
    \item \textbf{Capability Routing:} Execution is dynamically routed based on context
    \item \textbf{Capability Graph Execution:} Complex tasks decompose into capability graphs
\end{enumerate}

The execution layer functions as the action interface between AI cognition and the real-world capability ecosystem. This boundary decouples agents from providers, improves application portability, and supports ecosystem-scale growth.

More broadly, the proposed layer is analogous to operating-system and RPC boundaries in earlier computing eras: it provides a stable execution contract across heterogeneous providers while preserving higher-layer policy autonomy. As the ecosystem matures, this boundary is expected to become a foundational component of AI-native software infrastructure.
